% preamble.tex -- typesetting for Antifascist Intelligence.
% Hand-edited. This is the design surface: change type, measure, and styling
% here. Nothing regenerates this file.
%
% The document class is deliberately NOT declared here. It is in book.tex,
% above the \input that pulls this file in (D-170). Overleaf picks a project's
% main file by searching the .tex files for that declaration, so while it lived
% here Overleaf compiled preamble.tex, which has no \begin{document}, and
% aborted with "no legal \end found". The class options travel with the master
% for that reason. Do not move them back.
%
% The declaration is spelled out nowhere in this file, comment included, in
% case that search does not skip comments -- which is not checkable from here.

\usepackage{fontspec}
% The Ukrainian edition (manuscript/uk/book.tex defines \bookukrainian before it
% inputs this file) needs Cyrillic, which Pagella lacks. P052 is URW's Palladio,
% the face Pagella was made from, in the version that carries Cyrillic; it is a
% system font here (urw-base35), not part of TeX Live.
\ifdefined\bookukrainian
  \setmainfont{P052}
\else
  \setmainfont{TeX Gyre Pagella}          % a Palatino, matching the HTML path
\fi
\usepackage{microtype}
% The 2in side margins (D-774) leave a measure of about 4.3in, where justified
% 11pt text otherwise overruns the margin some 200 times; emergency stretch lets
% TeX loosen a line instead.
\setlength{\emergencystretch}{3em}
\usepackage[a4paper,vmargin=1in,hmargin=2in]{geometry}   % narrow measure, law-review style (D-774)
\usepackage{fancyhdr}
\usepackage{tcolorbox}
\usepackage{longtable,array}   % the condensed text's 18 tables, with long cells (D-621)
% The tables keep the 6.27in measure they were drawn for and stand an inch into
% each side margin (D-774): at the 4.3in text measure their narrowest columns
% fall under an inch and single words overrun them. Every table's column widths
% are fractions of \linewidth, so widening \linewidth inside the table and
% pulling it out by \LTleft and \LTright is the whole change.
\usepackage{etoolbox}
\AtBeginEnvironment{longtable}{%
  \setlength{\linewidth}{\dimexpr\textwidth+2in\relax}%
  \setlength{\LTleft}{-1in}\setlength{\LTright}{-1in}}
\tcbuselibrary{breakable}   % the only library the box style needs
\usepackage[english]{babel}
% TeX Live carries Ukrainian for babel as a locale (.ini), not an option.
\ifdefined\bookukrainian\babelprovide[import,main]{ukrainian}\fi
\usepackage{textcomp}
\usepackage{todonotes}
% refs.bib uses \euro, which is not a base LaTeX macro. Everything else it
% uses (\url \S \i \emph \c \v \textsection \L) is standard.
\providecommand{\euro}{\texteuro}
\usepackage[hidelinks]{hyperref}   % before biblatex, per its own guidance
% Citations are Chicago-style notes (D-774): every \autocite becomes a footnote,
% the first citation of a source in each chapter in full and later ones short,
% with no ibid. Chicago's own counts: up to three names in a note, more as the
% first and et al.; up to ten in the bibliography, more as the first seven and
% et al. biblatex-chicago is not installed here, so this is biblatex's verbose
% style with the adjustments below.
% Under tex4ht (the HTML proof) a short citation's link back to the bibliography
% breaks the footnote it sits in -- "Missing } inserted" in \@footnotetext -- so
% the HTML sets citations without links. The PDF keeps them.
\ifdefined\HCode\PassOptionsToPackage{hyperref=false}{biblatex}\fi
\usepackage[backend=biber,style=verbose,citereset=chapter,sorting=nyt,
  maxcitenames=3,mincitenames=1,maxbibnames=10,minbibnames=7]{biblatex}
% Let long URLs in the References break after digits and letters, not only at / and .:
% the OHCHR docstore links are one unbreakable string and ran off the page.
\setcounter{biburlnumpenalty}{100}
\setcounter{biburlucpenalty}{100}
\setcounter{biburllcpenalty}{100}
\addbibresource{../finishing/refs.bib}
\makeatletter
% A note prints in the footnote, at the end of the citation, after the URL and
% before any addendum, not where the standard drivers put it, between the title
% and the publisher (D-773). The bibliography carries no notes and no addenda
% (D-774).
\newcommand*{\af@savenote}{%
  \undef\af@note
  \iffieldundef{note}{}{\savefield{note}{\af@note}\clearfield{note}}}
\AtEveryCitekey{\af@savenote}
\AtEveryBibitem{\undef\af@note\clearfield{note}\clearfield{addendum}}
\renewbibmacro*{addendum+pubstate}{%
  \ifdef\af@note{\printtext{\af@note}\newunit\newblock}{}%
  \printfield{addendum}%
  \newunit\newblock
  \printfield{pubstate}}
\makeatother
% Chicago adjustments to the verbose style: a journal article has no "In:" before
% the journal; no access dates; and in the bibliography the first name of each
% entry is inverted for alphabetizing, the rest left in natural order.
\renewbibmacro*{in:}{%
  \ifentrytype{article}{}{\printtext{\bibstring{in}\intitlepunct}}}
\renewbibmacro*{url+urldate}{\usebibmacro{url}}
% The standard styles have no driver for court cases or statutes and fell back to
% one that dropped the court and the reporter citation ("Kareem v. Haspel.
% Judgment. Jan. 15, 2021."). A case prints as a report does -- title, type,
% number, court, date -- and a statute as a manual does, with its issuing body.
\DeclareBibliographyAlias{jurisdiction}{report}
\DeclareBibliographyAlias{legislation}{manual}
% A case or rule cited in reporter form carries its year in that citation --
% "323 U.S. 214 (1944)", "90 Fed. Reg. 8139 (Jan. 24, 2025)" -- and the date
% printed after it said the year twice (D-778). Where howpublished holds a
% parenthetical ending in a year, the date is not printed; it stays in the data
% for sorting. "(as amended through 2018)" is a revision, not the date, and is
% left alone.
\DeclareSourcemap{\maps[datatype=bibtex]{\map{
  \step[fieldsource=howpublished,
        match=\regexp{\s\((?![^()]*amended)[^()]*\d{4}\)}, final]
  \step[fieldset=keywords, fieldvalue={af:yearincitation}]}}}
\newcommand*{\afdropdate}{%
  \ifkeyword{af:yearincitation}
    {\clearfield{year}\clearfield{month}\clearfield{day}}{}}
\AtEveryCitekey{\afdropdate}
\AtEveryBibitem{\afdropdate}
\AtBeginBibliography{%
  \DeclareNameAlias{author}{family-given/given-family}%
  \DeclareNameAlias{editor}{family-given/given-family}%
  \DeclareNameAlias{translator}{family-given/given-family}}

% ---------------------------------------------------------------------------
% Draft apparatus: the DRAFT watermark and the status footer (D-319).
%
% One switch governs all of it. Set \draftmodefalse to ship a clean book;
% nothing else in this file has to change, and the two generated counts below
% go unread rather than stale.
\newif\ifdraftmode \draftmodetrue

% The two numbers the footer prints, and the list of entries a human has checked
% (one \refcheckedkey{KEY} per entry), which the citation highlight below reads.
% GENERATED by finishing/tools/refs_ledger.py from finishing/refs-ledger.tsv and
% verified current by check_all.sh, so a page cannot print a count or a check
% the ledger does not carry.
\newcommand{\refcheckedkey}[1]{\csdef{refok@#1}{}}
%% draft-status.tex -- GENERATED by finishing/tools/refs_ledger.py from
%% finishing/refs-ledger.tsv. Do not edit by hand; edits here are overwritten.
%%
%% The footer on every page prints these two numbers, and the draft highlight
%% leaves unmarked every citation of a key listed below by \refcheckedkey. They
%% are regenerated from the ledger and checked by check_all.sh, so the page cannot
%% claim a count or a check the ledger does not carry.
\newcommand{\refschecked}{8}
\newcommand{\refstotal}{612}
\refcheckedkey{abraham2024lavender}
\refcheckedkey{livingstone2026maven}
\refcheckedkey{Kaplan2026MinnProsecutors}
\refcheckedkey{Wikipedia2026ReneeGood}
\refcheckedkey{rickandmorty2023}
\refcheckedkey{gomezugarte2025uncertainty}
\refcheckedkey{acosta2026gaza}
\refcheckedkey{economist2025gaza}


% The build timestamp, MM/DD/YYYY HH:MM, computed by TeX as the job starts.
% Deliberately NOT written into draft-status.tex: a generated file that changed
% on every build would make that file's --check meaningless, and the check is
% the whole reason the footer can be trusted. Pure TeX arithmetic rather than
% \directlua so that the Overleaf round trip (D-169) still compiles if Overleaf
% picks an engine that is not LuaTeX.
\makeatletter
\newcount\ds@hh \newcount\ds@mm \newcount\ds@tmp
\newcommand{\ds@pad}[1]{\ifnum#1<10 0\fi\number#1}
\AtBeginDocument{%
  \ds@hh=\time \divide\ds@hh by 60
  \ds@mm=\time \ds@tmp=\ds@hh \multiply\ds@tmp by 60 \advance\ds@mm by -\ds@tmp
  \edef\buildstamp{%
    \ds@pad{\month}/\ds@pad{\day}/\number\year\space\ds@pad{\ds@hh}:\ds@pad{\ds@mm}}%
}
\makeatother

% The separator is a middot rather than the hyphen the line was sketched with:
% the book's own dashes are all spaced em dashes and a run of those in a
% four-part footer reads as heavily as the prose it sits under.
\newcommand{\draftstatusline}{%
  PREPRINT \textperiodcentered\ WORKING MANUSCRIPT \textperiodcentered\
  Rendered \buildstamp\ \textperiodcentered\
  \refschecked/\refstotal\ bibliographic entries human-checked}

\ifdraftmode
  \usepackage{xcolor}          % loaded by tcolorbox above; named here because
                               % \draftstatusline's grey depends on it
  \usepackage{draftwatermark}
  % Angle, size and position are the package's own defaults -- 45 degrees,
  % 0.25\paperwidth, centred -- and are left unset rather than restated. Only
  % the lightness moves: the default 0.8 grey is dark enough to read through on
  % a page of 11pt Pagella, and 0.88 stayed legible as a watermark without
  % competing with the prose it sits behind; 0.92, an 8% tint, is the author's
  % lighter setting (D-775).
  \SetWatermarkText{DRAFT}
  \SetWatermarkLightness{0.92}

  % Unchecked citations (D-641). Every in-text citation of an entry the ledger
  % does not mark checked is set on an orange highlight; a checked entry's
  % citations are set plain. The test is per key, inside biblatex's per-key
  % `cite' macro, so in \autocites{a}{b} each key carries its own state and the
  % brackets and pre/postnotes around them stay unmarked. To clear one, a human
  % marks it: finishing/tools/refs_ledger.py --mark KEY --by INITIALS (or
  % --toggle), then rebuild.
  %
  % Three renderings of the one mark. Under tex4ht (the HTML proof) a span the
  % page's stylesheet colours, since tex4ht can't carry lua-ul's highlight.
  % Under LuaTeX, lua-ul's \highLight, which breaks across lines and keeps the
  % citation's link working. Under any other engine -- an Overleaf compile on
  % XeLaTeX, say (D-169) -- orange text, because lua-ul needs LuaTeX.
  \ifdefined\HCode
    \newcommand{\refunchecked}[1]{%
      \HCode{<span class="ref-unchecked">}#1\HCode{</span>}}
  \else\ifdefined\directlua
    \usepackage{luacolor,lua-ul}
    \newcommand{\refunchecked}[1]{\highLight[orange!45]{#1}}
  \else
    \newcommand{\refunchecked}[1]{\textcolor{orange!80!black}{#1}}
  \fi\fi
  \letbibmacro{draft:cite}{cite}
  \renewbibmacro*{cite}{%
    \ifcsdef{refok@\thefield{entrykey}}
      {\usebibmacro{draft:cite}}
      {\refunchecked{\usebibmacro{draft:cite}}}}
\fi

% Under tex4ht, \appendix redefines \@chapter so that an appendix chapter no
% longer sets \@currentlabel, and a \label after it takes the number of the last
% section before it: the HTML proof printed Appendix A as "34", D as "C.3" and G
% as "F.6", with links to those sections. Setting the label after each appendix
% chapter restores the letter and the link (D-777). The PDF is unaffected.
\ifdefined\HCode
  \makeatletter
  \AddToHook{cmd/appendix/after}{%
    \let\af@appchapter\@chapter
    \def\@chapter[#1]#2{%
      \af@appchapter[#1]{#2}\protected@edef\@currentlabel{\thechapter}}}
  \makeatother
\fi

% Under tex4ht, each citation leaves its key in an HTML comment, invisible on
% the page. render_markdown.py reads the footnotes of the HTML build into the
% Markdown and matches each to the citation in the sources it came from by these
% keys, so a footnote can't land against the wrong citation (D-776).
\ifdefined\HCode
  \letbibmacro{md:cite}{cite}
  \renewbibmacro*{cite}{\HCode{<!--cite:\thefield{entrykey}-->}\usebibmacro{md:cite}}
\fi

\setcounter{secnumdepth}{2}             % deepest heading in the book is x.y.z
% The contents list chapters and sections and stops there. At depth 2 it ran
% pages 2-5 of a 196-page book, three full pages and an 11-line stub, and the 64
% subsections were most of it. What the depth was buying is a lookup path for the
% 79 cross-references that point at an x.y.z -- 35 percent of the 228 -- and the
% running head below now carries that instead, on the page the reader is already
% on. Cutting the depth without the head first would have left those 79 with no
% lookup at all, the class being oneside (D-260).
\setcounter{tocdepth}{1}

% Run-in head: a bold lead-in line above a paragraph, used 169 times.
\newcommand{\runin}[1]{\par\medskip\noindent\textbf{#1}\par\nopagebreak\smallskip}

% A one-line standing note under a heading, saying which of the preface's
% five claims the text beneath it rests on. Small italic, so it reads as a
% note and not as the first paragraph.
\newcommand{\standing}[1]{\noindent{\small\itshape #1}\par\medskip}

% Breakable by default; a box that reads as one document, like the birth
% record in chapter 31, opens with [unbreakable].
% In the Ukrainian edition the boxes don't break across pages. Under LuaTeX a
% hyperlink that falls on a breakable box's page break stops the build
% ("'endlink' ended up in different nesting level than 'startlink'"); the new
% face moved a citation in the box in 6.5 onto the break. The English book has
% the same exposure, but no link falls on a break there now.
\tcbset{esbreak/.style={breakable}}
\ifdefined\bookukrainian\tcbset{esbreak/.style={}}\fi
\newtcolorbox{esbox}[1][]{
  colback=black!3, colframe=black!35, boxrule=0.4pt, arc=1pt,
  left=8pt, right=8pt, top=8pt, bottom=8pt, esbreak,
  fontupper=\small,
  before upper={\setlength{\parskip}{0.5\baselineskip}%
               \setlength{\parindent}{0pt}},
  #1
}
\newcommand{\boxtitle}[1]{{\bfseries #1}\par\smallskip}

% Front and back matter are \chapter*, which does not set \@currentlabel, so a
% plain \label after one captures whatever number was last set -- silently
% wrong. This pins the value a \ref to the heading yields, which is the
% ordinal of the heading and not a number the page shows.
\makeatletter
\newcommand{\unnumberedlabel}[2]{\begingroup\def\@currentlabel{#2}\label{#1}\endgroup}
\makeatother

% A part epigraph sits on the part's title page, under the title. book.cls
% ends \part with \newpage, so a quote placed after it lands alone on the next
% page. When \partepigraph{...} immediately follows \part{...} (no blank line
% between), \@endpart takes it and sets it before the page break. Used
% anywhere else, or by an engine that doesn't run \@endpart (tex4ht), it
% prints as an ordinary quote, so the text is never lost.
\makeatletter
\newcommand{\partepigraph}[1]{\begin{quote}#1\end{quote}}
\let\af@endpart\@endpart
\long\def\af@partepigraph\partepigraph#1{%
  \vskip 3em{\normalsize\normalfont\begin{quote}#1\end{quote}}\af@endpart}
\def\@endpart{\@ifnextchar\partepigraph{\af@partepigraph}{\af@endpart}}
\makeatother

% \chapter* also sets no running mark, so the head left standing by the last
% numbered chapter persists over the back matter -- every page of the glossary
% read "CHAPTER 12. CONCLUSION AND OUTLOOK" (Q-055). Called once per starred
% chapter, after the heading.
\newcommand{\backmattermark}[1]{\markboth{\MakeUppercase{#1}}{\MakeUppercase{#1}}}

% The running head carries the deepest heading in force, so a reader who follows
% a cross-reference to an x.y.z can find it from the page they are on. The class
% is oneside, where book.cls drops \sectionmark entirely and puts the chapter on
% every page -- all 196 of them read "CHAPTER n. TITLE" and told a reader nothing
% they did not know. \subsectionmark is defined because \@sect calls it and the
% standard page styles \@gobble it, which is why the subsection never reached a
% head. Small caps rather than the class's uppercase: an x.y.z title set in full
% capitals is long and hard to read, and \backmattermark's own uppercase marks
% still set as capitals here, so the two agree (D-260).
\pagestyle{fancy}
\fancyhf{}
\fancyhead[L]{\small\scshape\rightmark}
\fancyhead[R]{\small\thepage}
\ifdraftmode\fancyfoot[C]{\scriptsize\color{black!55}\draftstatusline}\fi
\renewcommand{\headrulewidth}{0pt}
\renewcommand{\chaptermark}[1]{\markboth{\MakeUppercase{#1}}{\thechapter.\ #1}}
\renewcommand{\sectionmark}[1]{\markright{\thesection.\ #1}}
\renewcommand{\subsectionmark}[1]{\markright{\thesubsection.\ #1}}
% A chapter's opening page is \thispagestyle{plain}; it carries the folio alone.
\fancypagestyle{plain}{%
  \fancyhf{}%
  \fancyhead[R]{\small\thepage}%
  \ifdraftmode\fancyfoot[C]{\scriptsize\color{black!55}\draftstatusline}\fi%
  \renewcommand{\headrulewidth}{0pt}%
}

% The cover. book.tex calls \bookcover first of all, before \frontmatter, so
% the cover art (cover_cropped.png at the repository root, also the README's
% image) is the book's first page. In the PDF it is drawn behind an empty page,
% centred and as large as the page allows, with no footer and no watermark; the
% page is numbered A so that its PDF label and link target can't collide with
% the title page's i. In the HTML proof it is an ordinary image, which
% html_single_file.py embeds in the page.
\usepackage{graphicx}
\usepackage{eso-pic}
\newcommand{\bookcoverfile}{../cover_cropped.png}
\newcommand{\bookcover}{%
  \ifdefined\HCode
    \begin{center}\includegraphics[width=\linewidth]{\bookcoverfile}\end{center}%
    \clearpage
  \else
    \pagenumbering{Alph}%
    \ifdraftmode\DraftwatermarkOptions{stamp=false}\fi
    \AddToShipoutPictureBG*{\AtPageCenter{\makebox[0pt]{\raisebox{-0.5\height}{%
      \includegraphics[width=\paperwidth,height=\paperheight,keepaspectratio]%
        {\bookcoverfile}}}}}%
    \thispagestyle{empty}\null\clearpage
    \ifdraftmode\DraftwatermarkOptions{stamp=true}\fi
  \fi
}

% The title page. book.tex supplies \booktitlemain, \booksubtitle and
% \bookauthor, and calls \booktitlepage as the first thing after \frontmatter.
% The book had no title page at all until 2026-08-25 (D-073); it opened on the
% table of contents, so its own title appeared nowhere in the PDF.
% The notice under the title. The Ukrainian edition supplies its own before it
% inputs this file.
\providecommand{\booklegal}{%
     Copyright \textcopyright{} 2026 \bookauthor.\\
     Licensed under CC BY-NC-ND 4.0.\\[0.4\baselineskip]
     Third-party material quoted in this work, including the epigraphs,\\
     remains the property of its rights holders and is reproduced under\\
     fair use. The license above does not extend to it.}
\newcommand{\booktitlepage}{%
  \begin{titlepage}
    \centering
    \vspace*{0.28\textheight}
    {\fontsize{30}{36}\selectfont \booktitlemain\par}
    \vspace{1.2\baselineskip}
    {\Large\itshape \booksubtitle\par}
    \vspace{3\baselineskip}
    {\large \bookauthor\par}
    \vfill
    {\footnotesize\booklegal\par}
    \ifdraftmode
      \vspace{1.2\baselineskip}
      {\footnotesize\color{black!55}\draftstatusline\par}
    \fi
    \vspace{0.06\textheight}
  \end{titlepage}%
}
